\titre{}
\theme{ensemble}
\auteur{Nathan Scheinmann}
\niveau{1M}
\source{musy}
\type{serie}
\piments{2}
\pts{}
\annee{2425}

\contenu{
\tcblower
	On donne trois sous-intervalles de $\mathbb{R}: I=\interval[open right]{-3}{4}, J=\interval[open right]{-2}{0}$ et $K=\interval[open left]{-5}{3}$.

Donner à l'aide d'intervalles : $I \cap J, I \cup K, I \cap K, I \setminus K$ et $K \setminus I$.
}
\correction{
\tcblower
\begin{tasks}(2)
		\task $I\cap J=\interval[open right]{-3}{4}\cap \interval[open right]{-2}{0}=\interval[open right]{-2}{0}$ 
		\task $I\cup K=\interval[open right]{-3}{4}\cup \interval[open left]{-5}{3}=\interval[open]{-5}{4}$
		\task $I\cap K=\interval[open right]{-3}{4}\cap \interval[open left]{-5}{3} = \interval[]{-3}{3}$
		\task $I\setminus K=\interval[open right]{-3}{4}\setminus \interval[open left]{-5}{3}=\interval[open]{3}{4}$
		\task $K\setminus I=\interval[open left]{-5}{3}\setminus \interval[open right]{-3}{4}=\interval[open]{-5}{-3}$
	\end{tasks}
}

